\section{Conclusion}
\label{sec:conclusion}

\textsc{TAG} first attenuates unreliable instruction--response
supervision and then ranks candidates using a training-adaptive
difficulty--uncertainty carrier with bounded representation-anchor
modulation. Under covariance-drift and eigengap conditions, a local
Davis--Kahan analysis bounds one-refresh change of the recomputed
principal direction at the optimiser scale. This result concerns the
anchor direction only, while the factor
\(1+\lambda\widetilde a_i^{(t)}\in[1,1{+}\lambda]\) separately
limits its direct effect on the carrier; neither statement guarantees
stability of the full ranking. Empirically, the resulting selector improves the
main-setting average over static and policy-based dynamic baselines,
shows robust average performance in additional backbone and
dataset-transfer checks, and reduces end-to-end dynamic-selection cost
by about a third. These results suggest that representation geometry
recomputed from the current model state can provide an interpretable,
low-cost auxiliary signal for adaptive post-training data selection,
without requiring a learned selection policy.

